\section{A uniform bound through the gaps between Fourier levels}
We prove a lower bound for every sample size using a thin middle band
of Fourier modes whose covariance trace is small. Ordinary Lipschitz
contraction controls this band, while Lewis-weight sampling controls
the lower degrees.

\begin{lemma}[A fixed-threshold sampling bound]\label{lem:multiband-process}
Let $r<u<d$, and put $\Delta=\tau_r-\tau_u$. Let
$\mathcal B_{r,u}$ be the span of Walsh characters of degrees $r+1$
through $u$. Suppose
\[
 \eta^2\ge C\frac{D_r\log(2n)}n
\]
for a sufficiently large universal constant $C$. For deterministic
$R,T\ge0$, the class
\[
 \mathcal H=\{-f(p+b):p\in V_r,\ \mathbb E|p|\le R,
                 \ b\in\mathcal B_{r,u},\ \|b\|_K\le\sqrt n\}
\]
satisfies
\begin{equation}\label{eq:multiband-process}
 \mathbb E_S\sup_{h\in\mathcal H}
   |\mathbb P_n(h-T)_+-\mathbb E_x(h-T)_+|
 \le4(R\eta+\sqrt\Delta).
\end{equation}
Here $\mathbb P_n g=n^{-1}\sum_i g(X_i)$.
\end{lemma}
\begin{proof}
For the full-cube matrix with rows $f(x)(\chi_U(x))_{|U|\le r}$,
the Lewis weights are $D_r/M$. The linear-process moment bound of
Mai, Musco, and Rao \cite[Theorem 3]{mai2021coresets}, with its failure
parameter set to $1/2$, therefore implies
\[
 \mathbb E_{S,\epsilon}\sup_{p\in V_r,\,\mathbb E|p|\le R}
 \left|\frac1n\sum_i\epsilon_i f(X_i)p(X_i)\right|\le R\eta.
\]
The normalization follows directly: their sampled matrix rescales each
sampled row by $M/n$, whereas $\|p\|_1=M^{-1}\|A\beta\|_1$.
Their moment bound implies the displayed expectation bound by Jensen.

The covariance-weighted feature vector of the band has squared norm
$\Delta$ at every input. Consequently
\[
 \mathbb E_{S,\epsilon}\sup_{\|b\|_K\le\sqrt n}
 \left|\frac1n\sum_i\epsilon_i f(X_i)b(X_i)\right|
 \le\frac{\sqrt n}{n}\sqrt{n\Delta}=\sqrt\Delta.
\]
The map $z\mapsto(z-T)_+$ is $1$-Lipschitz and vanishes at zero.
Symmetrization, followed by the scalar contraction principle, bounds
the left-hand side of \eqref{eq:multiband-process} by four times the
linear Rademacher width of $\mathcal H$. The triangle inequality and
the preceding two estimates prove the claim. Both parameter balls and
the threshold are fixed before sampling; no positive homogeneity in
$T$ or relative bound on an unbounded parameter class is asserted.
\end{proof}

